% TOP_PROBLEM: {"id": 20000968, "title": "A q-Eulerian last-site law for cycle IDLA", "category_id": 2, "category": {"id": 2, "name": "combinatorics", "display_name": "Combinatorics", "description": "Counting problems, graph theory, discrete structures.", "slug": "combinatorics", "order_index": 2, "created_at": "2026-07-31T15:26:25.671Z"}, "status": "partially_solved", "problem_number": "AIM-COMBINATORICS-0093", "set_id": 14}
% FIRST_POSTED: 2026-10-01
% ENTRY_KIND: statement_only
% UnsolvedMath ID 20000968; source catalogue code AIM-COMBINATORICS-0093
% !TeX program = lualatex
% Definition and sources only; this file is not an LLM solution attempt.
\documentclass[11pt,a4paper]{article}
\usepackage[margin=25mm]{geometry}
\usepackage{fontspec}
\setmainfont{lmroman10-regular.otf}[BoldFont=lmroman10-bold.otf,
 ItalicFont=lmroman10-italic.otf,BoldItalicFont=lmroman10-bolditalic.otf]
\usepackage{amsmath,amssymb,mathtools}
\usepackage{unicode-math}
\setmathfont{latinmodern-math.otf}
\AtBeginDocument{\let\setminus\smallsetminus}
\usepackage{xurl,hyperref}
\hypersetup{colorlinks=true,urlcolor=blue}
\setcounter{secnumdepth}{0}
\setlength{\parindent}{0pt}
\setlength{\parskip}{5pt}
\setlength{\emergencystretch}{3em}
\begin{document}
\section*{A q-Eulerian last-site law for cycle IDLA}\label{20000968}
{\small UnsolvedMath ID 20000968 · AIM-COMBINATORICS-0093 · Combinatorics · AIM Workshop Problem Lists}
\subsection{Definitions and mathematical statement}
Let $G=(V,E)$ be a finite connected undirected graph with at least two vertices, a distinguished root $q$, and positive edge conductances $w_{uv}=w_{vu}$.
A random walk at $u$ moves to a neighbor $v$ with probability $w_{uv}/\sum_{z\sim u}w_{uz}$.
Start with the occupied set $A_1=\{q\}$.
For $j=1,\ldots,|V|-1$, launch a new independent walk from $q$, stop it when it first visits $V\setminus A_j$, and add that stopping vertex to form $A_{j+1}$.
Connectivity ensures that each stopping time is finite almost surely.
This is the finite explorers game, or internal diffusion-limited aggregation.

Let $L$ be the unique vertex in $A_{|V|}\setminus A_{|V|-1}$, and write
\[
 p_{G,q,w}(v)=\mathbb P(L=v),\qquad v\in V.
\]
The workshop asks for a useful description of this last-occupied-vertex distribution beyond the unweighted cycle, including higher-dimensional graphs, circulant graphs, and weighted cycles.
A circulant graph here has vertex set $\mathbb Z/m\mathbb Z$ and edges whose allowed differences form a fixed symmetric generating set.
For weighted cycles the positive conductances are part of the input and may break rotational symmetry.
Seek exact formulas or effective recurrences that reveal how the distribution depends on the graph and weights, and asymptotic descriptions for natural growing families.
The terminal vertex is defined by successive explorer particles; it is not the last newly visited vertex of one uninterrupted random walk.
\subsection{Short English statement}
Describe the distribution of the last site occupied by successive explorers on graphs, especially circulants, weighted cycles, and higher-dimensional examples.
\subsection{Sources}
\begin{itemize}
\item[S1] ULAM.AI and the UnsolvedMath Contributors, supplied problems.json, record 20000968 (AIM-COMBINATORICS-0093), AIM Workshop Problem Lists. Catalogue identity and source transcription; CC BY 4.0.
\url{https://huggingface.co/datasets/ulamai/UnsolvedMath}.
\item[S2] American Institute of Mathematics, Generalizations of chip-firing and the critical group, item 3. Original workshop question underlying AIM-COMBINATORICS-0093.
\url{https://aimath.org/pastworkshops/chipfiringproblems.pdf}.
\end{itemize}
\end{document}
